\section{Physical motivation}

The isomer shift $\IS$ measures the shift of the center of the Mössbauer
pattern with respect to a reference, usually $\alpha$-Fe. It is related to the
electron density at the nucleus and, therefore, to the oxidation state,
coordination and covalency of the Fe atom. In ordered materials a single $\IS$
is usually fitted per phase. However, in glasses, silicates, amorphous
materials, ferritins, nanoparticles or mixtures of local environments, a
continuous distribution of shifts may exist.

Fitbauer implements two extensions:
\begin{itemize}
  \item a one-dimensional distribution $P(\IS)$;
  \item two-dimensional distributions that combine $\IS$ with another quantity,
  especially $P(\IS,\QS)$ and $P(\BHF,\IS)$.
\end{itemize}
The $P(\IS,\QS)$ case is particularly useful for broad paramagnetic doublets,
where the center of the doublet and its separation may vary simultaneously.

\begin{figure}[htbp]\centering
  \includegraphics[width=0.95\textwidth]{apx_is}
  \caption{Isomer shift distribution. \emph{Left}: a
  $P(\IS)$ with two populations separated in $\delta$ —for example Fe$^{3+}$
  ($\delta\!\approx\!0{,}35$~mm/s) and Fe$^{2+}$ ($\delta\!\approx\!1{,}05$~mm/s)—
  reflects two oxidation states or environments. \emph{Right}: the resulting
  spectrum is the superposition of the patterns of each $\delta$; when the
  environments overlap, the distribution separates them better than a rigid
  discrete fit.}
  \label{fig:apx_is}
\end{figure}

\section{One-dimensional distribution $P(\IS)$}

For a grid of shifts $\IS_i$, the 1D model is
\begin{equation}
  T(v_k)=b_0+b_1v_k-\sum_i P_i\,S(v_k;\IS_i,\QS,\BHF,\Gamma),
\end{equation}
where $S$ is the unit-depth absorption pattern. If $\BHF=0$, the pattern reduces
in practice to a doublet or singlet depending on the value of $\QS$. The weights
satisfy
\begin{equation}
  P_i\ge0.
\end{equation}
The regularized problem is written as
\begin{equation}
  \Phi(\bm p)=\|W(\bm y-X\bm p)\|^2+\alpha\|D^2\bm p\|^2,
\end{equation}
where $D^2$ is the second-difference operator. The normalized probability is
computed as
\begin{equation}
  p(\IS_i)=\frac{P_i}{\int P(\IS)\,\dd\IS}.
\end{equation}

\section{Two-dimensional distribution $P(\IS,\QS)$}

For a grid $\IS_i$, $Q_j$, the model is
\begin{equation}
  T(v_k)=b_0+b_1v_k-\sum_{ij}P_{ij}\,
  S(v_k;\IS_i,Q_j,\BHF,\Gamma).
\end{equation}
If one wishes to describe a paramagnetic population, $\BHF=0$ is fixed and each
point $(\IS_i,Q_j)$ generates a doublet centered at $\IS_i$ with separation
$Q_j$:
\begin{align}
  v_{1,ij}&=\IS_i-\frac{Q_j}{2},\\
  v_{2,ij}&=\IS_i+\frac{Q_j}{2}.
\end{align}

The 2D regularization penalizes curvature in both directions:
\begin{equation}
  \Phi(\bm P)=\|W(\bm y-X\bm z)\|^2
  +\alpha_{\IS}\|D_{\IS}^2P\|^2
  +\alpha_Q\|D_Q^2P\|^2,
  \qquad P_{ij}\ge0.
\end{equation}
In augmented form:
\begin{equation}
  \begin{bmatrix}
    WX\\
    \sqrt{\alpha_{\IS}}\,D_{\IS}^2\\
    \sqrt{\alpha_Q}\,D_Q^2
  \end{bmatrix}\bm z\simeq
  \begin{bmatrix}
    W\bm y\\
    \bm 0\\
    \bm 0
  \end{bmatrix}.
\end{equation}

\section{Interpretation of $P(\IS,\QS)$}

The marginals are
\begin{align}
  P_{\IS}(\IS_i)&=\int P(\IS_i,Q)\,\dd Q,\\
  P_Q(Q_j)&=\int P(\IS,Q_j)\,\dd\IS.
\end{align}
A tilted cloud in the map may indicate that the sites with higher $\IS$ also
have a higher or lower $\QS$. This may suggest families of local environments,
for example variation in valence, coordination or polyhedron distortion.

Fitbauer computes descriptive diagnostics of the map:
\begin{align}
  \langle\IS\rangle,\quad \sigma_{\IS},\quad
  \langle Q\rangle,\quad \sigma_Q,\quad
  \rho_{\IS Q}.
\end{align}
The apparent correlation is
\begin{equation}
  \rho_{\IS Q}=\frac{\langle(\IS-\langle\IS\rangle)(Q-\langle Q\rangle)\rangle}
  {\sigma_{\IS}\sigma_Q}.
\end{equation}
These values describe the fitted distribution, but on their own they do not
prove that the distribution is unique or physical.

\section{Distribution $P(\BHF,\IS)$}

$P(\BHF,\IS)$ is also implemented:
\begin{equation}
  T(v_k)=b_0+b_1v_k-\sum_{ij}P_{ij}\,
  S(v_k;B_i,\QS,\IS_j,\Gamma).
\end{equation}
It can be useful in magnetic phases where the hyperfine field and the isomer
shift vary jointly, for example due to a size distribution, partial oxidation or
surface environments. However, it is especially sensitive to velocity
calibration errors: a global shift of the scale can look like an $\IS$
distribution.

\section{Identifiability risks}

\textbf{Important warning:} distributions that contain $\IS$ are very sensitive
to calibration, folding and background. In a doublet,
\begin{equation}
  v_1=\IS-\frac{\QS}{2},\qquad v_2=\IS+\frac{\QS}{2},
\end{equation}
so that changes in $\IS$ shift the center and changes in $\QS$ modify the
separation. With broad lines or noise, many $(\IS,\QS)$ combinations can produce
similar spectra. In addition, $\IS$ correlates with:
\begin{itemize}
  \item the folding point;
  \item the $V_{\max}$ calibration;
  \item the isomeric reference;
  \item the width $\Gamma$;
  \item the slope/background;
  \item omitted discrete components.
\end{itemize}
Therefore, a $P(\IS)$ or $P(\IS,\QS)$ distribution that is visually smooth and
has a good residual may lack physical meaning if it is compensating for an
instrumental error or an incomplete model.

\section{Good practices}

Before interpreting a distribution involving $\IS$ it is recommended to:
\begin{enumerate}
  \item fix or carefully verify the velocity calibration;
  \item use a known isomeric reference;
  \item check the folding point with the residual diagnostic;
  \item compare with discrete models and with the 1D $P(\QS)$;
  \item vary $\alpha_{\IS}$ and $\alpha_Q$ over several orders of magnitude;
  \item repeat with a different number of bins;
  \item avoid publishing fine details that disappear when the regularization is
  changed;
  \item use external information: chemistry, XRD, TEM, magnetometry or the
  literature.
\end{enumerate}
For publication, it is advisable to show that the $P(\IS,\QS)$ map is stable and
that a simpler model does not explain the data satisfactorily.
